% ============================================================
%  ch02b  §2.2 函数域上的整体算术朗兰兹
%  原文对应：ICM-arXiv.tex 937--1009 行
%  编号：定理 2.2.1；注 2.2.2；公式 (2.5)；猜想 2.2.3；
%        定理 2.2.4；推论 2.2.5
% ============================================================

\subsection{函数域上的整体算术朗兰兹}\label{S: Global Langlands function fields}

接下来转向算术朗兰兹对应的整体方面。如开头所述，其经典表述粗略地预测（不可约）Galois 表示的集合与（尖点）自守表示的集合之间存在自然的对应。与局部情形一样，除 $\GL_n$ 情形（这是 \cite{Lafforgue-GLn} 的定理）之外，这样的表述并不容易精确化。另一方面，\S\ref{Global GLC} 中的整体几何朗兰兹猜想以及去范畴化／迹的哲学提示：整体算术朗兰兹可以、而且多半应当表述为两个向量空间之间的同构，这两个空间分别来自 Galois 一侧与自守一侧。本小节在整体函数域情形表述这样一个猜想。

设 $F=\bF_q(X)$ 是 $\bF_q$ 上几何连通光滑射影曲线 $X$ 的函数域。以 $\eta=\Spec F$ 记 $X$ 的生成点，$\overline\eta$ 记 $\eta$ 上的一个几何点。以 $|X|$ 记 $X$ 的闭点集。对 $v\in |X|$，以 $\mO_v$ 记 $X$ 在 $v$ 处的完全局部环，$F_v$ 为其分式域。令 $\bO_F=\prod_{v\in|X|}\mO_v$ 为整 adèle 环，$\bA_F=\prod'_{v\in|X|}F_v$ 为 adèle 环。对有限非空的位置集 $Q$，以 $W_{F,Q}$ 记 $F$ 的在 $Q$ 外非分歧的 Weil 群。

设 $G$ 是 $F$ 上的连通约化群。与局部情形类似，表述整体猜想的第一步是 \cite{Zhu-Conjecture} 中的下述定理。

\begin{thm}
设 $\ell\nmid 2p$。把 $\bZ_\ell$-代数 $A$ 送到满足 $d\circ\rho=(\cycl^{-1},\pr)$ 的（强）连续同态 $\rho: W_{F,Q}\to \CG(A)$ 所构成空间的预叠，可由一个导出概形 $\Loc_{\CG,Q}^\Box$ 表示；该概形是若干导出仿射概形的无交并，每个在 $\bZ_\ell$ 上平坦、有限型。若 $Q\neq\emptyset$，则 $\Loc_{\CG,Q}^\Box$ 拟光滑。
\end{thm}

于是定义整体朗兰兹参数叠为 $\Loc_{\CG,Q}=\Loc_{\CG,Q}^\Box/\hat{G}$。与局部情形类似（见例 \ref{Ex: trace of Satake}），对 $\hat{G}_\La$ 在有限投射 $\La$-模上的表示，把它看作 $\bB\hat{G}_\La$ 上的向量丛，以 $\widetilde V$ 记其拉回至 $\Loc_{\CG,Q}\otimes\La$。若 $V$ 是 $(\CG)^S$ 的某个表示沿对角嵌入 $\hat{G}\to (\CG)^S$ 的限制，则 $\widetilde V$ 上有自然的（强）连续 $W_{F,Q}^S$-作用（见 \cite[Sect. 2.4]{Zhu-Conjecture}）。对 $F$ 的位置 $v$，以 $\Loc_{\CG,v}$ 记 $G_{F_v}$ 的局部朗兰兹参数叠。令
\[
\res: \Loc_{\CG,Q}\to \prod_{v\in Q}\Loc_{\CG,v}
\]
为由整体参数限制到局部参数所得的映射（它由 $W_{F_v}\to W_{F,Q}$ 诱导）。下文将考虑沿此映射对 $\prod_{v\in Q}\Loc_{\CG,v}$ 上凝聚层取 $!$-拉回。

\begin{rmk}
(1) 事实上，当 $Q=\emptyset$ 时 $\Loc_{\CG,Q}$ 的定义需要略作修改。

(2) 与局部情形不同，$\Loc_{\CG,Q}$ 一般带有非平凡的导出结构（见 \cite[Rem. 3.4.5]{Zhu-Conjecture}）。以 ${}^{cl}\Loc_{\CG,Q}$ 记其底层的经典叠。

(3) \cite{AGKRRV} 给出了 $\Loc_{\CG,Q}\otimes\bQ_\ell$ 的另一种定义。
\end{rmk}

接着转到自守一侧。为简单起见，本小节设 $G$ 在 $\bF_q$ 上分裂。固定一个水平，即开紧子群 $K\subset G(\bO_F)$。设 $Q$ 是由那些满足 $K_v\neq G(\mO_v)$ 的 $v$ 构成的位置集。对有限集 $S$，以 $\Sht_K(G)_{(X-Q)^S}$ 记 $(X-Q)^S$ 上由 $X$ 上带 $K$-水平结构的 $S$-leg（legs）落在 $X-Q$ 中的 $G$-shtuka 的模空间所给出的 ind-Deligne--Mumford 叠。（该模空间的基本构造与性质例如见 \cite{Lafforgue-Automorphic-to-Galois}。）它沿对角映射 $\overline\eta\to (X-Q)\xrightarrow{\Delta}(X-Q)^S$ 的基变换记作 $\Sht_K(G)_{\Delta(\overline\eta)}$。

对 $(\CG)^S$ 在有限投射 $\La$-模上的每个表示 $V$，几何 Satake \eqref{factorizable Sat}（把 $D$ 换成 $X-Q$，并允许 $\La=\bZ_\ell$）给出 $\Sht_K(G)_{(X-Q)^S}$ 上的反常层 $\on{Sat}_S(V)$。以 $C_c\bigl(\Sht_K(G)_{\Delta(\overline\eta)},\Sat_S(V)\bigr)$ 记 $\Sht_K(G)_{\Delta(\overline\eta)}$ 以 $\Sat_S(V)$ 为系数的总的紧支撑上同调（的上链复形）。它带有（强）连续的 $W_{F,Q}^S$-作用（导出层次上这种作用的构造见 \cite{HRS-Kunneth}，其基础是 \cite{Xue-Finiteness, Xue-Smoothness}），还带有相应的整体（导出）Hecke 代数（系数在 $\La$ 中）
\begin{equation}\label{E: derived hecke global}
H_{K,\La}=\Big(R\End \big(c\on{-ind}_K^{G(\bA_F)}(\La)\big)\Big)^{\op}
\end{equation}
的作用。例如，若 $V=\mathbf{1}$ 是平凡表示，则（在 $G$ 分裂的假设下）
\[C_c\bigl(\Sht_K(G)_{\Delta(\overline\eta)}, \on{Sat}_{\{1\}}(\mathbf{1})\bigr)= C_c(G(F)\backslash G(\bA)/K,\La).\]
这里 $G(F)\backslash G(\bA)/K$ 看作 $\overline\eta$ 上的离散 DM 叠，而 $C_c(G(F)\backslash G(\bA)/K,\La)$ 记其紧支撑上同调。当 $\La=\bQ_\ell$ 时，这就是 $G(F)\backslash G(\bA)/K$ 上紧支撑函数构成的空间。

我们将固定 $G$ 的一个钉 $(B,T,e)$ 以及非退化特征 $\psi: F\backslash \bA\to \bZ_\ell[\mu_{p}]^\times$，它按猜想 \ref{C: LLC} 对每个 $v\in Q$ 给出猜想性的等价 $\bL_{v}$。特别地，与 $K_v\subset G(F_v)$ 相对应，有 $\Loc_{\CG,v}$ 上的猜想凝聚层 $\frakA_{K_v}$（见 \eqref{E: sheaf A}）。

\begin{conj}\label{C: gen coh}
存在自然的 $(W_{F,Q}^S\times H_{K,\La})$-等变同构
\[
R\Gamma\bigl(\on{Loc}_{\CG,Q}\otimes\La, \widetilde{V}\otimes  \res^!(\boxtimes_{v\in Q} \frakA_{K_v})\bigr)\cong C_c\bigl(\on{Sht}_K(G)_{\Delta(\overline\eta)}, \on{Sat}_S(V)\bigr).
\]
\end{conj}

该猜想更一般的形式（允许“广义水平结构”）以及在各种特殊情形下的例子，见 \cite[Sect. 4.7]{Zhu-Conjecture}。这个猜想可以看作函数域上整体朗兰兹对应的精确形式：它给出了把 Galois 表示与自守表示对应起来的精确配方。（例如，V. Lafforgue 的巡游算子就编码在此同构之中，见下文。）此外，这样的同构符合 Arthur--Kottwitz 重数公式，同时把该公式推广到整系数层次，从而与自守提升理论相关联。

该猜想最令人信服的证据是下述定理 \cite{Lafforgue-Zhu-Elliptic, Zhu-Conjecture}，它正是 Drinfeld 对 Lafforgue 构造所作解释（在启发式层次上）所提示的。

\begin{thm}\label{T: gen LZ}
对每个 $i$，存在 ${}^{cl}\Loc_{\CG,Q}\otimes\bQ_\ell$ 上的拟凝聚层 $\frakA^i_K$，带有 $H_{K,\bQ_\ell}$ 的作用，使得对 $(\CG)^S$ 的每个有限维 $\bQ_\ell$-表示 $V$，都有自然的 $(W_{F,Q}^S\times H_{K,\bQ_\ell})$-等变同构
\begin{equation*}\label{E: coh Sht}
\Gamma\big({}^{cl}\Loc_{\CG,Q}\otimes\bQ_\ell,\widetilde V\otimes \frakA^i_K\big)\cong H^i_c\bigl(\Sht_K(G)_{\Delta(\overline\eta)},\Sat_S(V)\bigr).
\end{equation*}
\end{thm}

需要指出，\cite{Lafforgue-Zhu-Elliptic, Zhu-Conjecture} 实际上对任意 $G$ 都证明了这个定理。此外，当 $K$ 处处超特殊时，\cite{AGKRRV} 证明 \eqref{E: coh Sht} 在导出层次上成立。

同构 \eqref{E: coh Sht} 诱导出 $\Gamma\big({}^{cl}\Loc_{\CG,Q}\otimes\bQ_\ell,\mO\big)$ 在右端的作用。这正是 V. Lafforgue 的巡游算子所给出的作用，它诱导出右端（特别是 $C_c(G(F)\backslash G(\bA_F)/K, \bQ_\ell)$）按半单朗兰兹参数的分解。如 \cite{Lafforgue-Zhu-Elliptic} 所说明，在椭圆朗兰兹参数上，这一同构与 Arthur--Kottwitz 重数公式密切相关。当 $G=\GL_n$ 时，它给出下述推论，推广了 \cite{Lafforgue-GLn}。

\begin{cor}
设 $\pi$ 是 $\GL_n$ 的尖点自守表示，其相伴的不可约 Galois 表示 $\rho_\pi: W_{F,Q}\to \GL_n(\La)$ 取值于某个有限扩张 $\La/\bQ_\ell$，并以 $\frakm_\pi$ 记 $\Gamma\big({}^{cl}\Loc_{\CG,Q}\otimes\La,\mO\big)$ 中相应的极大理想。则存在 $(W_{F,Q}^S\times H_K)$-等变同构
\[
H^*_c\big(\Sht_K(G)_{\Delta(\overline\eta)}, \Sat_S(V)\big)/\frakm_\pi\cong  V_{\rho_\pi}\otimes \pi^K.
\]
特别地，左端只集中在上同调次数零。
\end{cor}
